\documentclass[11pt,a4paper]{article}
\usepackage[T1]{fontenc}
\usepackage[utf8]{inputenc}
\usepackage{amsmath,amssymb,amsthm}
\usepackage[margin=1in]{geometry}
\usepackage[colorlinks=true,linkcolor=blue,urlcolor=blue]{hyperref}
\theoremstyle{plain}
\newtheorem{theorem}{Theorem}
\newtheorem{lemma}[theorem]{Lemma}
\newtheorem{proposition}[theorem]{Proposition}
\newtheorem{corollary}[theorem]{Corollary}
\theoremstyle{definition}
\newtheorem{remark}[theorem]{Remark}

\title{The empirical closed form for OEIS A185474, proved}
\author{Adrian Perez Fontelles\\ \small Independent researcher}
\date{1 September 2026}
\begin{document}
\maketitle

\begin{abstract}
OEIS A185474 carries an empirical closed form for $a(n)$ --- not a recurrence relating
consecutive terms but an explicit formula. It is true, and it is decidable rather than
empirical. The entry counts arrays subject to a condition on a bounded window of consecutive lines, so the lines of an array are the
vertices of a finite digraph, an array is a walk in it, and $a(n)$ is a walk count on
$S=314213$ vertices. Any function of the form $\sum_j p_j(n)\lambda_j^{\,n}$ satisfies the
monic linear recurrence whose characteristic polynomial is
$\prod_j(x-\lambda_j)^{\deg p_j+1}$; here that polynomial is $q(x)=(x-1)^{20}$, of degree
$20$. Two things are then checked exactly: that the walk count satisfies the same
recurrence from some index on, which the Cayley--Hamilton theorem reduces to a finite
computation, and that the two sequences agree at $20$ consecutive indices past that
point. A monic recurrence determines every later term from its predecessors, so agreement
there is agreement for good.
\end{abstract}

\noindent\small 2020 Mathematics Subject Classification. 05A15, 11B37, 15A18.\normalsize

\section{The conjecture}

OEIS A185474 is ``Number of (n+2) X 8 0..2 arrays with each 3 X 3 subblock having rows and columns in lexicographically nondecreasing order.''. It has offset $1$ and begins
\[
169052, 1433903, 8297059, 37961900, 146203201, 493061605, 1497314456, 4179700035,\ \dots
\]
The entry states, as a conjecture contributed by its author and never marked settled:
\begin{quote}
Empirical: a(n) = (1/212666259456000)*n\^{}19 + (67/83147710464000)*n\^{}18 + (12559/194011324416000)*n\^{}17 + (12833/3923023104000)*n\^{}16 + (333169/2853107712000)*n\^{}15 + (196437919/62768369664000)*n\^{}14 + (1153058911/17118646272000)*n\^{}13 + (1339332949/1034643456000)*n\^{}12 + (28852162981/1207084032000)*n\^{}11 + (350922659539/877879296000)*n\^{}10 + (4735468557467/877879296000)*n\^{}9 + (11975048873639/219469824000)*n\^{}8 + (131882024236729/329204736000)*n\^{}7 + (7064899018536199/3362591232000)*n\^{}6 + (1862906706349061/237758976000)*n\^{}5 + (13955757767803297/653837184000)*n\^{}4 + (26671169520943/623750400)*n\^{}3 + (851721864110749/15437822400)*n\^{}2 + (4258571767339/116396280)*n + 2793.
\end{quote}
As of the ``Last modified'' line on the live entry (Jan 18 2022 00:35:56, revision 7) this is still
recorded as empirical, and nothing on the entry records it as proved. Write
\[
f(n)\;=\;\frac{n^{19}}{212666259456000} + \frac{67 n^{18}}{83147710464000} + \frac{12559 n^{17}}{194011324416000} + \frac{12833 n^{16}}{3923023104000} + \frac{333169 n^{15}}{2853107712000} + \frac{196437919 n^{14}}{62768369664000} + \frac{1153058911 n^{13}}{17118646272000} + \frac{1339332949 n^{12}}{1034643456000} + \frac{28852162981 n^{11}}{1207084032000} + \frac{350922659539 n^{10}}{877879296000} + \frac{4735468557467 n^{9}}{877879296000} + \frac{11975048873639 n^{8}}{219469824000} + \frac{131882024236729 n^{7}}{329204736000} + \frac{7064899018536199 n^{6}}{3362591232000} + \frac{1862906706349061 n^{5}}{237758976000} + \frac{13955757767803297 n^{4}}{653837184000} + \frac{26671169520943 n^{3}}{623750400} + \frac{851721864110749 n^{2}}{15437822400} + \frac{4258571767339 n}{116396280} + 2793.
\]

\section{The count is a walk count}

The entry's defining condition is local: it is a condition on a bounded window of consecutive lines. Take as vertices the
admissible configurations of such a window and put an edge from $u$ to $v$ when $v$ may
follow $u$, which the condition decides by inspection. An admissible array is then exactly a
walk, so with $M$ the adjacency matrix of that digraph on $S=314213$ vertices, and $u,w$ the
vectors recording which windows may start and end an array,
\[
a(n)\;=\;u^{\!\top}M^{\,g(n)}w
\]
for the linear function $g$ that the entry's shape supplies. In particular $a$ is
$C$-finite: it satisfies some monic integer linear recurrence, though not necessarily a short
one.

\section{A closed form is a recurrence}

\begin{lemma}\label{lem:ann}
Let $f(m)=\sum_{j}p_j(m)\lambda_j^{\,m}$ with the $\lambda_j$ distinct and the $p_j$
polynomials, and put $q(x)=\prod_j(x-\lambda_j)^{\deg p_j+1}$. Then $q$ applied to the shift
operator annihilates $f$: writing $q(x)=x^{r}-\sum_{i=1}^{r}c_ix^{r-i}$,
\[
f(m)\;=\;\sum_{i=1}^{r}c_i\,f(m-i)\qquad\text{for every }m.
\]
\end{lemma}

\begin{proof}
The shift operator $E$ acts on $m\mapsto\lambda^{m}p(m)$ as
$(E-\lambda)\bigl(\lambda^{m}p(m)\bigr)=\lambda^{m+1}\bigl(p(m+1)-p(m)\bigr)$, and
$p(m+1)-p(m)$ has degree one less than $p$. So $(E-\lambda)^{\deg p+1}$ kills
$\lambda^{m}p(m)$, and the factors of $q$ commute.
\end{proof}

Here $q(x)=(x-1)^{20}$, of degree $20$; the closed form is a polynomial in $n$, so this is $(x-1)^{20}$, and the recurrence it encodes is
\begin{equation}\label{eq:rec}
a(n)\;=\;20\,a(n-1) - 190\,a(n-2) + 1140\,a(n-3) - 4845\,a(n-4) + 15504\,a(n-5) - 38760\,a(n-6) + 77520\,a(n-7) - 125970\,a(n-8) + 167960\,a(n-9) - 184756\,a(n-10) + 167960\,a(n-11) - 125970\,a(n-12) + 77520\,a(n-13) - 38760\,a(n-14) + 15504\,a(n-15) - 4845\,a(n-16) + 1140\,a(n-17) - 190\,a(n-18) + 20\,a(n-19) - a(n-20).
\end{equation}

\section{The proof}

\begin{lemma}\label{lem:crit}
With $q$ as above, $a$ satisfies \eqref{eq:rec} for all large $n$ if and only if
$u^{\!\top}M^{\,j}q(M)w=0$ for every $j\ge0$, and it is enough to check $j<S$.
\end{lemma}

\begin{proof}
Substituting the walk expression, the residual $a(n)-\sum_ic_ia(n-i)$ equals
$u^{\!\top}M^{\,g(n)-r}q(M)w$ up to the fixed linear reparametrisation $g$. For the bound,
$M^{S}$ is an integer combination of $I,M,\dots,M^{S-1}$ by Cayley--Hamilton, so the values
$u^{\!\top}M^{j}q(M)w$ satisfy the monic recurrence given by the characteristic polynomial of
$M$; $S$ consecutive zeros therefore force all later ones, and the last nonzero value pins the
threshold exactly.
\end{proof}

\begin{theorem}
$a(n)=f(n)$ for every $n\ge1$.
\end{theorem}

\begin{proof}
The vector $w'=q(M)w$ was formed by $20$ matrix--vector products in exact integer
arithmetic, and $u^{\!\top}M^{j}w'$ evaluated for $j=0,1,\dots$, again exactly. Those integers
vanish from the point corresponding to $n=21$ onwards, and the run was continued until the
working vector was identically zero, which settles every larger $j$ at once. So $a$ satisfies
\eqref{eq:rec} for every $n>20$. By Lemma~\ref{lem:ann} $f$ satisfies \eqref{eq:rec}
identically. Both were then evaluated in exact arithmetic at the $20$ consecutive indices
$n=21,\dots,40$ and agree at each. A monic recurrence of order $20$
determines $a(m)$ from $a(m-1),\dots,a(m-20)$, and for every $m\ge41$ both
$m>20$ and $m-20\ge21$ hold, so induction on $m$ gives $a(m)=f(m)$ for every
$m\ge21$.
Below $n=21$ the recurrence argument gives nothing, since \eqref{eq:rec} is only
established for $a$ from $n=21$ on. The remaining indices $n=1,\dots,20$ were
therefore checked one at a time, directly against the walk count in exact integer arithmetic,
and agree.
\end{proof}

The range is tight in the sense that was checked: the closed form holds from the entry's first term onwards.

\section{Verification}

Four checks, each independent of the algebra above.

First, the walk model reproduces every one of the $17$ terms the entry publishes,
exactly, in integer arithmetic. This is what ties the model to the entry: the model is built
by reading the entry's English, and a misreading gives a different digraph and different
counts.

Second, the closed form was evaluated in exact rational arithmetic --- not floating point ---
at every published term from $n=1$ on, and agrees with the entry's own DATA at each.

Third, the annihilation test of Lemma~\ref{lem:crit} was carried out exactly, and the model
and the closed form were then compared directly at sixty consecutive indices beyond
$n=1$, far past where the proof already settles the matter.

Fourth, the closed form was deliberately perturbed --- by a constant and by a term linear in
$n$ --- and each perturbed form was rejected by the same comparison, so the test is not one
that anything would pass.

\begin{thebibliography}{9}
\bibitem{oeis} The OEIS Foundation, \emph{The On-Line Encyclopedia of Integer
Sequences}, \url{https://oeis.org}, 2026. Sequence A185474.
\bibitem{stanley} R.~P.~Stanley, \emph{Enumerative Combinatorics}, Volume 1, 2nd ed.,
Cambridge University Press, 2011. (Transfer-matrix method, Section 4.7; $C$-finite sequences,
Section 4.1.)
\bibitem{kp} M.~Kauers and P.~Paule, \emph{The Concrete Tetrahedron}, Springer, 2011.
(Closed forms and linear recurrences with constant coefficients.)
\bibitem{horn} R.~A.~Horn and C.~R.~Johnson, \emph{Matrix Analysis}, 2nd ed., Cambridge
University Press, 2013.
\end{thebibliography}

\end{document}
